\documentclass{beamer}
\usetheme{Darmstadt}
\usepackage{comment}
\usepackage[utf8]{inputenc}    % Para escribir acentos y caracteres especiales
\usepackage{graphicx}          % Para insertar imágenes
\usepackage{booktabs}          % Para tablas más bonitas
\usepackage{enumitem}
\usepackage{amsmath}           % Para escribir ecuaciones
\usepackage{tikz}              % Para gráficos vectoriales
\usetikzlibrary{calc}          % Para coordenadas calculadas con TikZ
\usepackage{pgfplots}          % Para gráficos matemáticos
\pgfplotsset{compat=1.18}      % Evita errores por compatibilidad
\usepackage{hyperref}
\usepackage{pgf}
\usepackage{ragged2e}

\newif\ifpresentacion
%\presentaciontrue  % Cambiar a \presentacionfalse para PDF sin animaciones
\presentacionfalse % Cambiar a \presentaciontrue para PDF con  animaciones

\usepackage{ifthen}
\newcommand{\pausa}{\ifpresentacion\pause\fi}
% #region Título
\title{Ayudantía 3 \\ Valoración de Futuros y
Forwards \\ \large\textit{Instrumentos Derivados}}
\author{
  \texorpdfstring{
    \textbf{Profesor:} Francisco Rantul \\[0.3em]
    \textbf{Ayudante:} Mateo Canales
  }{Profesor: Francisco Rantul, Ayudante: Mateo Canales}
}
\subject{Instrumentos Derivados}
\institute{Universidad Diego Portales}
\date{16 De Abril, 2025}
% #endregion 

\begin{document}

% Portada
\begin{frame}
    \titlepage
    \vfill
    \centering
    \includegraphics[width=2.3118cm]{../imagenes/logo.png}
  \end{frame}

% #region Formateo de números 
\newcommand{\entero}[1]{\pgfmathprintnumber[fixed, precision=0]{#1}}
\newcommand{\decimal}[1]{\pgfmathprintnumber[fixed, precision=2]{#1}}
\newcommand{\decimalx}[1]{\pgfmathprintnumber[fixed, precision=3]{#1}}
\newcommand{\decimalxx}[1]{\pgfmathprintnumber[fixed, precision=4]{#1}}
\newcommand{\porcentaje}[1]{%
  \pgfmathsetmacro{\temp}{#1*100}%
  \pgfmathprintnumber[fixed, precision=2]{\temp}\%%
  }
\newcommand{\dinero}[1]{%
  \$\,\pgfmathprintnumber[fixed, precision=3]{#1}
  }
\pgfkeys{/pgf/number format/.cd,1000 sep={\,}} % usa espacio fino como separador

\newcommand{\miles}[1]{\pgfmathprintnumber[fixed, precision=0]{#1}}
\newcommand{\forward}{\ensuremath{F_0 =  S_0 \cdot{e^{r \cdot T}}}}
\newcommand{\forwarddividendo}{\ensuremath{F_0 =  \left(S_0-I \right) \cdot{e^{r \cdot T}}}}
\newcommand{\forwardalmacenamiento}{\ensuremath{F_0 =  \left(S_0+U \right) \cdot{e^{r \cdot T}}}}
\newcommand{\capcontinuacero}{ F = P \cdot e^{-rT}}
\newcommand{\forwardmonedas}{\ensuremath{F_0 =  S_0 \cdot{e^{\left(r-r_r\right) \cdot T}}}}

% #endregion

% #region Variables pregunta 4

\newcommand{\plata}{18.8}
\newcommand{\almacenamiento}{0.4}
\newcommand{\tlr}{0.04}
\newcommand{\pdolardos}{708}
\newcommand{\tclpdos}{0.0264}
\newcommand{\tusddos}{0.01}
\newcommand{\espflujouno}{\ensuremath{P = E \left(\sum_{i=1}^{n} Flujo_i \right) \cdot{e^{-r \cdot T}}}}
\newcommand{\tcuatroa}{\frac{9}{12}}
\pgfmathsetmacro{\tcuatroac}{9/12}
\newcommand{\tcuatroab}{\frac{6}{12}}
\pgfmathsetmacro{\tcuatroabc}{6/12}
\newcommand{\tcuatroaa}{\frac{3}{12}}
\pgfmathsetmacro{\tcuatroaac}{3/12}
\pgfmathsetmacro{\flujocuatroa}{\almacenamiento*exp{-\tlr*\tcuatroaac}}
\pgfmathsetmacro{\flujocuatrob}{\almacenamiento*exp{-\tlr*\tcuatroabc}}
\pgfmathsetmacro{\flujocuatroc}{\almacenamiento*exp{-\tlr*\tcuatroac}}
\pgfmathsetmacro{\flujocuatro}{\flujocuatroa+\flujocuatrob+\flujocuatroc}
\pgfmathsetmacro{\valoresecero}{\plata+\flujocuatro }
\pgfmathsetmacro{\expcuatro}{\tlr*\tcuatroac }
\pgfmathsetmacro{\finalcuatro}{\valoresecero * exp{\expcuatro}}

% #endregion

\section{Pregunta 4}

\begin{frame}{Pregunta 4}
  \justify
  Suponga que el precio spot del commodity de plata es actualmente igual a \dinero{\plata} dólares  
  la onza. Los costos de almacenamiento son iguales a \dinero{\almacenamiento}por año la onza, pagaderos por
  trimestres vencidos. La estructura de tasas de interés es plana con una tasa cero libre 
  de riesgo del \porcentaje{\tlr} anual compuesto continuo.
  \begin{enumerate}[label=\textbf{\alph*)}]
  \item   Se le pide calcular cual debería ser el precio de futuros de plata, con entrega a 9 meses plazo.
  \item 	Explique que ocurre si el precio de futuros de plata con entrega a 9 meses,
  efectivamente observado en el mercado es de \$19,8 dólares la onza.
  \item 	Asuma que existe una baja demanda por el commodity de plata, por lo tanto, las empresas están 
  almacenando altos niveles de inventarios. Explique qué ocurre con el rendimiento de conveniencia.
  
\end{enumerate}
\end{frame}

\subsection{Parte a)}

\begin{frame}{Pregunta 4 parte a)}
  \justify

Suponga que el precio spot del commodity de plata es actualmente igual a \dinero{\plata} dólares  
la onza. Los costos de almacenamiento son iguales a \dinero{\almacenamiento} por año la onza, pagaderos por
trimestres vencidos. La estructura de tasas de interés es plana con una tasa cero libre 
de riesgo del \porcentaje{\tlr} anual compuesto continuo.\\
\vspace{.5em}
\textbf{a)} Se le pide calcular cual debería ser el precio de futuros de plata, con entrega a 9 meses plazo.

\end{frame}

\begin{frame}{Desarrollo a)}
  \justify

  Datos: \(S_0=\plata\), \(almacenamiento=\almacenamiento\), \(r=\tlr\), \(T=\tcuatroa\).\\
  \pausa
  Anteriormente usamos una fórmula para dividendos, y que si lo pensamos de cierta manera, 
  el almacenamiento es lo mismo que un dividendo pero el flujo es contrario.\\
  \pausa
  \vspace{.5em}
  \(\forwarddividendo\  \pausa \rightarrow\ \) \textcolor{blue}{\(\forwardalmacenamiento\)}\\
  \pausa
  Siendo U el valor presente del valor de almacenamiento. con la fórmula de valorización continua \footnote{ \(\capcontinuacero\)}.\\
  \pausa
  \(U=Alm_1 + Alm_2 + Alm_3\)\\
  \pausa
  \(U= \almacenamiento \cdot e^{-\tlr \cdot \tcuatroaa} + \almacenamiento \cdot e^{-\tlr \cdot \tcuatroab} + \almacenamiento \cdot e^{-\tlr \cdot \tcuatroa}\)\\
  \pausa
  \(U = \decimalx{\flujocuatroa} + \decimalx{\flujocuatrob} + \decimalx{\flujocuatroc}\)\\
  \pausa
  \fcolorbox{black}{green!10}{
  \(U = \decimalx{\flujocuatro}\)
  } 
\end{frame}

\begin{frame}{Desarrollo a) (continuación)}
  \Large
  Datos: \(S_0=\plata\), \(r=\tlr\), \(T=\tcuatroa\),  \(U = \decimalx{\flujocuatro}\).\\
  \textcolor{blue}{\(\forwardalmacenamiento\)}\\
  \pausa
  Reemplazando\\
  \(F_0= \left(\plata + \decimalx{\flujocuatro}\right) \cdot e^{\tlr \cdot \tcuatroa } \)\\
  \pausa
  \(F_0= \left(\decimalx{\valoresecero}\right) \cdot e^{\decimalx{\expcuatro}} \)\\
  \pausa
  \fcolorbox{black}{green!10}{
  \(F_0= \decimalx{\finalcuatro}\)
  }
\end{frame}


\subsection{Parte b)}

\begin{frame}{Pregunta 4 parte b)}
  \justify
  \large
Suponga que el precio spot del commodity de plata es actualmente igual a \dinero{\plata} dólares  
la onza. Los costos de almacenamiento son iguales a \dinero{\almacenamiento} por año la onza, pagaderos por
trimestres vencidos. La estructura de tasas de interés es plana con una tasa cero libre 
de riesgo del \porcentaje{\tlr} anual compuesto continuo.\\
\vspace{.5em}
\textbf{b)} Explique que ocurre si el precio de futuros de plata con entrega a 9 meses,
efectivamente observado en el mercado es de \$19,8 dólares la onza.

\end{frame}

\begin{frame}{Desarrollo b)}
\justify
\Large
\textbf{Respuesta:}  
Según la fórmula \(\forwardalmacenamiento\), el valor calculado debería ser igual al valor que está en el mercado,
por ende, es correcto asumir que existirian posibilidades de arbitraje.
\end{frame}

\subsection{Parte c)}

\begin{frame}{Pregunta 4 parte c)}
  \justify
  \large
  Suponga que el precio spot del commodity de plata es actualmente igual a \dinero{\plata} dólares  
  la onza. Los costos de almacenamiento son iguales a \dinero{\almacenamiento} por año la onza, pagaderos por
  trimestres vencidos. La estructura de tasas de interés es plana con una tasa cero libre 
  de riesgo del \porcentaje{\tlr} anual compuesto continuo.\\
  \vspace{.5em}
  \textbf{c)}	Asuma que existe una baja demanda por el commodity de plata, por lo tanto, 
  las empresas están almacenando altos niveles de inventarios. 
  Explique qué ocurre con el rendimiento de conveniencia. 
\end{frame}

\begin{frame}{Desarrollo c)}
  \justify
  \large
  Recordando, los rendimiento de conveniencia $(y)$ son los beneficios derivados de la tenencia del activo físico.
  Es decir, en algunos casos es recomendable mantener el activo como activo de consumo (materia prima) que como activo
  de inversión (futuro).\\
  \(F_0 \leq \left(S_0+U\right) \cdot e^{rT} \)\\
  \pausa
  Incluyendo el retorno de conveniencia:\\
  \(F_0 \cdot e^{yT}= \left(S_0+U\right) \cdot e^{rT} \)\\
  \pausa
  \(F_0 = \left(S_0+U\right) \cdot \frac{e^{rT}}{e^{yT}}  \)\\
  \pausa
  \(F_0 = \left(S_0+U\right) \cdot e^{\left(r-y\right)T} \)\\
\end{frame}
  
\begin{frame}{Desarrollo c) (continuación)}
  \justify
  \Large
  \(F_0 = \left(S_0+U\right) \cdot e^{\left(r-y\right)T} \)\\
 \vspace{1em} 
  Analizando ésta fórmula y el caso. La demanda cae, por ende el valor de tener el activo físico no es 
  tan alto (ya que puedo comprarlo en cualquier momento o no tiene plusvalía), es decir $y \rightarrow 0$, 
  aumentando el valor del forward.


\end{frame}
  
% #region Variables pregunta 5
\newcommand{\tcincounomeses}{6} 
\newcommand{\tcincouno}{\frac{6}{12}} 
\pgfmathsetmacro{\tcincounoc}{6/12}
\newcommand{\litros}{1060000}
\section{Pregunta 5}
\begin{frame}{Pregunta 5}
  \justify
  Suponga que LAN está preocupada por las variaciones del precio del petróleo durante los siguientes meses,
  específicamente los siguientes \textbf{\tcincounomeses\ meses}, ya que el combustible que ocupa la flota de aviones ocupa un importante 
  ítem de costo en el presupuesto de la firma, LAN sabe que necesitara \miles{\litros} litros de combustible en \textbf{\tcincounomeses\ meses}
  más para sus operaciones nacionales.\\
  El actual precio de combustible para aviones está en USD 65 el barril (159 litros). \\
  El gran problema que posee LAN es que el precio del combustible depende directamente del precio del Barril de Petróleo 
  en el mercado internacional el cual es actualmente de USD 48,6 y solo existen derivados sobre petróleo y no sobre combustible.
  
  \end{frame}
  
  \begin{comment}
  \begin{frame}{Pregunta 5}
  \begin{enumerate}[label=\textbf{\alph*)}]
  \small
  \item Calcule el precio teórico de un futuro sobre petróleo a 6 meses si existe un costo trimestral de almacenaje y 
  envió de USD 0,5 y la tasa libre de riesgo relevante es de 3\% anual (compuesta continua).
  \item Suponga que la desviación estándar de los cambios semestrales en los precios del combustible es USD 1.231 y 
  que la desviación estándar de los cambios semestrales en el precio del futuro del petróleo es de USD 1.285, y el coeficiente 
  de correlación entre estos cambios es de 0.94. Calcule el ratio de cobertura (hedge ratio) óptimo para un contrato a 6 meses 
  y explique cómo LAN podría ocupar el número calculado.
  \item Calcule cuantos contratos de petróleo necesitara LAN para la cobertura de combustible si cada contrato es de 150 barriles.
  \item Cuál sería el resultado de la estrategia de LAN si en 6 meses más el precio del petróleo fuera de USD 57,5 el barril y el de combustible de USD 75,5.
  \item Cuál sería el resultado de la estrategia de LAN si en 6 meses más el precio del petróleo fuera de USD 41,5 el barril y el de combustible de USD 58,1.
  \end{enumerate}
  \end{frame}


  \subsection{Parte a)}
  \begin{frame}{Pregunta 5 parte a}
  \textbf{a)} Calcule el precio teórico de un futuro sobre petróleo a 6 meses si existe un costo trimestral de almacenaje y 
  envió de USD 0,5 y la tasa libre de riesgo relevante es de 3\% anual (compuesta continua).

  \end{frame}

  \subsection{Parte b)}
  \begin{frame}{Pregunta 5 parte b}
  \textbf{b)} Suponga que la desviación estándar de los cambios semestrales en los precios del combustible es USD 1.231 y 
  que la desviación estándar de los cambios semestrales en el precio del futuro del petróleo es de USD 1.285, y el coeficiente 
  de correlación entre estos cambios es de 0.94. Calcule el ratio de cobertura (hedge ratio) óptimo para un contrato a 6 meses 
  y explique cómo LAN podría ocupar el número calculado.
  \end{frame}

  \subsection{Parte c)}
  \begin{frame}{Pregunta 5 parte c}
  \textbf{c)} Calcule cuantos contratos de petróleo necesitara LAN para la cobertura de combustible si cada contrato es de 150 barriles.
  \end{frame}

  \subsection{Parte d)}
  \begin{frame}{Pregunta 5 parte d}
  \textbf{d)}  Cuál sería el resultado de la estrategia de LAN si en 6 meses más el precio del petróleo fuera de USD 57,5 el barril y el de combustible de USD 75,5.
  \end{frame}

  \subsection{Parte e)}
  \begin{frame}{Pregunta 5 parte e}
  \textbf{e)} Cuál sería el resultado de la estrategia de LAN si en 6 meses más el precio del petróleo fuera de USD 41,5 el barril y el de combustible de USD 58,1.
  \end{frame}

  \section{Pregunta 6}
  \begin{frame}{Pregunta 6}
  El gerente de un fondo de inversión en acciones nacionales mantiene una cartera valorada en 18 millones de US\$ con un Beta de 
  1,25 en septiembre del 2022; el gerente está preocupado sobre el comportamiento del mercado durante los próximos once meses y 
  piensa en utilizar contratos de futuro a doce meses sobre el índice IPSA para cubrir el riesgo. El nivel del índice el 01/09/2022 
  es de 4.055 y un contrato de futuros es sobre 250 veces el índice. La tasa de interés libre de riesgo es de 3\% anual, la tasa de 
  rendimiento por dividendos sobre las acciones del índice es de 0.8\% anual y el tipo de cambio de dicho día fue de \$800  Se pide:
  \end{frame}

  \begin{frame}{Pregunta 6}
  \begin{enumerate}[label=\textbf{\alph*)}]
  \item Cuál es el precio teórico del futuro para el contrato de futuros de 12 meses.
  \item Qué posición en contratos de futuros debe tomar el gerente del Fondo, para realizar una cobertura óptima a lo largo de los próximos 12 meses.
  \item Calcule el efecto de su estrategia sobre el rendimiento del Fondo que administra el gerente, si el nivel del mercado en septiembre del 2023 del índice es 4.144.
  \item Calcule el efecto de su estrategia sobre el rendimiento del Fondo que administra el gerente, si el nivel del mercado en septiembre del 2023 del índice es 3.615.

  \end{enumerate}
  \end{frame}

  \subsection{Parte a)}

  \begin{frame}{Pregunta 6 parte a)}
  \textbf{a)} Cuál es el precio teórico del futuro para el contrato de futuros de 12 meses.
  \end{frame}

  \subsection{Parte b)}
  \begin{frame}{Pregunta 6 parte b)}
  \textbf{b)} Qué posición en contratos de futuros debe tomar el gerente del Fondo, para realizar una cobertura óptima a lo largo de los próximos 12 meses.
  \end{frame}

  \subsection{Parte c)}
  \begin{frame}{Pregunta 6 parte c)}
  \textbf{c)} Calcule el efecto de su estrategia sobre el rendimiento del Fondo que administra el gerente, si el nivel del mercado en septiembre del 2023 del índice es 4.144.
  \end{frame}

  \subsection{Parte d)}
  \begin{frame}{Pregunta 6 parte d)}
  \textbf{d)} Calcule el efecto de su estrategia sobre el rendimiento del Fondo que administra el gerente, si el nivel del mercado en septiembre del 2023 del índice es 3.615.
  \end{frame}
  
\end{comment}

\end{document}